\\documentclass[11pt]{article}
\\usepackage[T1]{fontenc}
\\usepackage[utf8]{inputenc}
\\usepackage{geometry}
\\usepackage{xcolor}
\\usepackage{fvextra}
\\usepackage{hyperref}
\\geometry{margin=1in}
\\DefineVerbatimEnvironment{CodeBlock}{Verbatim}{breaklines,breakanywhere,fontsize=\\small}

\begin{document}
\section*{core/spectral\\_polarity.py}
\textbf{Source Path}: \texttt{core/spectral\\_polarity.py}\par
\bigskip
\begin{CodeBlock}
"""
Track 1.5: Spectral Polarity — The Endogenous Compass

Performs PCA (via SVD) on the gradient matrix G formed by the 8 probe-pair
deltas (cls_per_bot). Extracts:
  - u_axis:  PC1, the direction of maximum semantic stress
  - EVR:     explained variance ratio (dipole validity gate)
  - probe_magnitudes:  signed projections of each probe onto u_axis
  - antagonism:  u_axis scaled by dominant singular value (directional force)

Multi-Scale Persistence (ASTER v3.2):
  - Sweep across bandwidths σ = [0.5, 1.0, 2.0] (or adaptive)
  - Compute u_axis(σ), EVR(σ) at each scale
  - Valid if EVR >= τ in at least min_persistent_scales (default 2)
  - Final u_axis = normalized mean over valid scales

Gauge fixing ensures consistent orientation across articles by projecting
a reference probe onto u_axis and flipping if negative.

References:
  ASTER Protocol v3.2, Section I (Track 1.5: Gradient Metrics / Compass)
"""

from dataclasses import dataclass, field
from typing import List, Optional, Tuple

import torch
import numpy as np


@dataclass
class SpectralPolarityConfig:
    """Configuration for spectral polarity decomposition."""
    # Index of the reference probe for gauge fixing (0 = first pair in framing_queries.yaml)
    reference_probe_idx: int = 0
    # EVR threshold for dipole validity (tau). Below this, the compass is unreliable.
    evr_threshold: float = 0.5
    # Multi-scale bandwidth values (sigma sweep)
    sigma_values: Tuple[float, ...] = (0.5, 1.0, 2.0)
    # Minimum number of scales that must pass EVR threshold for persistence
    min_persistent_scales: int = 2
    # Whether to use adaptive bandwidth (median heuristic) as one of the scales
    use_adaptive_sigma: bool = True
    # Whether to enable multi-scale mode (False = single-scale legacy mode)
    multi_scale_enabled: bool = True
    # Adaptive sigma sampling cap for median-heuristic pairwise distances.
    adaptive_sigma_max_samples: int = 1000
    # Lower bound clamp for adaptive sigma before any downstream scaling/division.
    adaptive_sigma_min: float = 1e-6
    # Floor for all sigma divisions in G / sigma conditioning paths.
    sigma_scale_epsilon: float = 1e-6
    # Shared numerical floor for variance/weight normalization.
    normalization_epsilon: float = 1e-12
    # Dynamic-K component threshold ratio (sigma_i > ratio * sigma_1).
    dynamic_k_threshold_ratio: float = 0.1
    # Explicit spectral weighting mode:
    # - "sigma": legacy sigma * delta^2
    # - "sigma_squared": sigma^2 * delta^2 (default)
    weighted_sigma_mode: str = "sigma_squared"
    # Explicit whitened weighting mode:
    # - "sigma": legacy (delta / sigma)^2
    # - "sigma_squared": (delta / sigma^2)^2
    whitened_sigma_mode: str = "sigma"
    # Floor for whitened sigma term before inverse weighting.
    whitened_sigma_epsilon: float = 1e-6


@dataclass
class SpectralPolarityResult:
    """Per-batch results from spectral polarity decomposition."""
    # [N, H] — first right singular vector (compass direction) per article
    u_axis: torch.Tensor
    # [N] — explained variance ratio of PC1 (aggregated across valid scales)
    evr: torch.Tensor
    # [N] — boolean mask: True where EVR >= threshold (valid dipole)
    dipole_valid: torch.Tensor
    # [N, 8] — signed projection of each probe gradient onto u_axis
    probe_magnitudes: torch.Tensor
    # [N, 8] — full singular value spectrum per article (from dominant scale)
    singular_values: torch.Tensor
    # [N, H] — directional antagonism vector = u_axis * S[0]
    antagonism: torch.Tensor
    # [N] — number of scales where EVR >= threshold (persistence count)
    n_persistent_scales: torch.Tensor
    # [N, n_scales] — EVR at each scale (for diagnostics)
    evr_per_scale: torch.Tensor
    # Dipole state string: "DIPOLE_STATE" or "BLINKER_STATE"
    dipole_state: str = "DIPOLE_STATE"
    # [N] — dynamic K: number of significant components per article
    dynamic_k: Optional[torch.Tensor] = None
    # [N, K, H] — top-K right singular vectors (u_basis) for weighted distance
    u_basis: Optional[torch.Tensor] = None


class SpectralPolarity:
    """
    Spectral decomposition of the probe-pair gradient matrix.

    Given G = [8, H] per article (the 8 mean-pool deltas from NLI extraction),
    computes the SVD to find the principal axis of semantic stress (u_axis),
    gauge-fixes for sign consistency, and produces a directional antagonism
    vector for downstream Track 5 fusion.

    Multi-Scale Persistence:
    Sweeps across multiple bandwidths to ensure the detected dipole is not
    a scale-dependent artifact. Only dipoles that persist across multiple
    scales are considered valid.
    """

    def __init__(self, config: Optional[SpectralPolarityConfig] = None):
        self.config = config or SpectralPolarityConfig()

    def _sanitize_sigma(self, sigma: float, min_sigma: float) -> float:
        """Clamp sigma to a finite positive floor for stable scaling."""
        sigma_value = float(sigma)
        if not np.isfinite(sigma_value):
            return float(min_sigma)
        return float(max(sigma_value, min_sigma))

    def _sigma_weight_term(self, sigma: torch.Tensor, mode: str) -> torch.Tensor:
        """Convert singular values into explicit sigma weighting semantics."""
        if mode == "sigma":
            return sigma
        if mode == "sigma_squared":
            return sigma ** 2
        raise ValueError(f"Unknown sigma weighting mode: {mode}")

    def _compute_adaptive_sigma(self, G: torch.Tensor) -> float:
        """
        Compute adaptive bandwidth using median heuristic.

        σ = median of pairwise distances between probe gradients.
        This adapts to the local density of the gradient manifold.
        """
        N, n_probes, H = G.shape
        # Flatten to [N*n_probes, H] for pairwise distance
        G_flat = G.reshape(-1, H)
        # Sample if too large (avoid O(n^2) explosion)
        max_samples = max(int(self.config.adaptive_sigma_max_samples), 1)
        if G_flat.shape[0] > max_samples:
            idx = torch.randperm(G_flat.shape[0])[:max_samples]
            G_sample = G_flat[idx]
        else:
            G_sample = G_flat
        # Pairwise L2 distances
        dists = torch.cdist(G_sample, G_sample, p=2)
        # Median of non-zero distances
        mask = dists > 0
        if mask.sum() > 0:
            median_dist = dists[mask].median().item()
        else:
            median_dist = 1.0
        return self._sanitize_sigma(median_dist, self.config.adaptive_sigma_min)

    def _compute_single_scale(
        self,
        G: torch.Tensor,
        sigma: float,
    ) -> Tuple[torch.Tensor, torch.Tensor, torch.Tensor]:
        """
        Compute SVD at a single bandwidth scale.

        Args:
            G: [N, 8, H] gradient matrix
            sigma: bandwidth scaling factor

        Returns:
            u_axis: [N, H] principal direction
            evr: [N] explained variance ratio
            S: [N, 8] singular values
        """
        # Scale the gradient matrix by bandwidth.
        safe_sigma = self._sanitize_sigma(sigma, self.config.sigma_scale_epsilon)
        G_scaled = torch.nan_to_num(
            G / safe_sigma, nan=0.0, posinf=0.0, neginf=0.0
        )

        # SVD: G = U S Vt
        U, S, Vt = torch.linalg.svd(G_scaled, full_matrices=False)

        # PC1: direction of maximum gradient variance
        u_axis = Vt[:, 0, :]  # [N, H]

        # EVR = S[0]^2 / sum(S^2)
        S_sq = S ** 2
        total_var = S_sq.sum(dim=1).clamp(min=self.config.normalization_epsilon)
        evr = S_sq[:, 0] / total_var

        return u_axis, evr, S

    def _gauge_fix(
        self,
        u_axis: torch.Tensor,
        G: torch.Tensor,
        ref_idx: int,
    ) -> torch.Tensor:
        """
        Apply gauge fixing to ensure consistent orientation.

        Projects reference probe gradient onto u_axis and flips if negative.
        """
        ref_gradient = G[:, ref_idx, :]  # [N, H]
        dot = (ref_gradient * u_axis).sum(dim=1)  # [N]
        sign = torch.sign(dot)
        # Handle zero dot products: default to +1
        sign = torch.where(sign == 0, torch.ones_like(sign), sign)
        return u_axis * sign.unsqueeze(1)

    def compute_batch(
        self,
        cls_per_bot_list: List[torch.Tensor],
    ) -> SpectralPolarityResult:
        """
        Compute spectral polarity for a batch of articles.

        Args:
            cls_per_bot_list: List of [8, H] tensors, one per article.
                Each row is the mean-pool delta (hypothesis_A - hypothesis_B)
                for one of the 8 framing probe pairs.

        Returns:
            SpectralPolarityResult with all per-article spectral diagnostics.
        """
        device = cls_per_bot_list[0].device
        dtype = cls_per_bot_list[0].dtype

        # Stack into [N, 8, H]
        G = torch.stack(cls_per_bot_list, dim=0).to(dtype=torch.float32)
        N, n_probes, H = G.shape

        if not self.config.multi_scale_enabled:
            # Legacy single-scale mode
            return self._compute_single_scale_legacy(G, dtype)

        # --- Multi-Scale Sigma Sweep ---
        sigma_list = list(self.config.sigma_values)

        # Add adaptive sigma if enabled
        if self.config.use_adaptive_sigma:
            adaptive_sigma = self._compute_adaptive_sigma(G)
            if adaptive_sigma not in sigma_list:
                sigma_list.append(adaptive_sigma)
        sigma_list = [
            self._sanitize_sigma(sigma, self.config.sigma_scale_epsilon)
            for sigma in sigma_list
        ]

        n_scales = len(sigma_list)

        # Storage for per-scale results
        u_axes_per_scale = []  # List of [N, H]
        evr_per_scale = []     # List of [N]
        S_per_scale = []       # List of [N, 8]
        master_u_axis = None

        for i, sigma in enumerate(sigma_list):
            u_axis_s, evr_s, S_s = self._compute_single_scale(G, sigma)
            # Align per-scale u-axis sign to avoid cancellation when aggregating.
            if master_u_axis is None:
                master_u_axis = u_axis_s.clone()
            if i > 0:
                sign_flip = (master_u_axis * u_axis_s).sum(dim=1) < 0
                u_axis_s = torch.where(sign_flip.unsqueeze(1), -u_axis_s, u_axis_s)
                # Maintain a running master direction so each scale aligns to the
                # accumulated axis, not only the first scale.
                master_u_axis = master_u_axis + u_axis_s
                # Preserve magnitude; do not L2-normalize away thermodynamic scale.
                master_u_axis = master_u_axis
            u_axes_per_scale.append(u_axis_s)
            evr_per_scale.append(evr_s)
            S_per_scale.append(S_s)

        # Stack for analysis: [N, n_scales, ...]
        evr_stacked = torch.stack(evr_per_scale, dim=1)  # [N, n_scales]
        u_axes_stacked = torch.stack(u_axes_per_scale, dim=1)  # [N, n_scales, H]

        # --- Persistence Check ---
        # Count how many scales pass EVR threshold per article
        scale_valid = evr_stacked >= self.config.evr_threshold  # [N, n_scales]
        n_persistent = scale_valid.sum(dim=1)  # [N]

        # Dipole is valid if it persists across min_persistent_scales
        dipole_valid = n_persistent >= self.config.min_persistent_scales  # [N]

        # --- Aggregate u_axis across valid scales ---
        # Mask invalid scales, compute weighted mean by EVR
        evr_masked = evr_stacked * scale_valid.float()  # Zero out invalid
        evr_weights = evr_masked / evr_masked.sum(dim=1, keepdim=True).clamp(
            min=self.config.normalization_epsilon
        )  # [N, n_scales]

        # Weighted sum of u_axes
        # [N, n_scales, H] * [N, n_scales, 1] -> sum over scales -> [N, H]
        u_axis_aggregated = (u_axes_stacked * evr_weights.unsqueeze(2)).sum(dim=1)

        # L2 normalize the aggregated direction
        # Preserve magnitude; do not L2-normalize away thermodynamic scale.
        u_axis_aggregated = u_axis_aggregated

        # For articles with no valid scales, fall back to scale with highest EVR
        no_valid_mask = n_persistent == 0
        if no_valid_mask.any():
            best_scale_idx = evr_stacked.argmax(dim=1)  # [N]
            for i in range(N):
                if no_valid_mask[i]:
                    u_axis_aggregated[i] = u_axes_stacked[i, best_scale_idx[i]]

        # --- Gauge Fixing on aggregated u_axis ---
        u_axis_fixed = self._gauge_fix(u_axis_aggregated, G, self.config.reference_probe_idx)

        # --- Aggregated EVR ---
        # Mean EVR across valid scales (or max EVR if none valid)
        evr_aggregated = torch.where(
            n_persistent > 0,
            (evr_stacked * scale_valid.float()).sum(dim=1) / n_persistent.clamp(min=1),
            evr_stacked.max(dim=1).values
        )

        # --- Probe Magnitudes ---
        probe_magnitudes = torch.bmm(
            G, u_axis_fixed.unsqueeze(2)
        ).squeeze(2)  # [N, 8]

        # --- Singular Values (from scale=1.0 or first valid) ---
        # Use sigma=1.0 as reference, or find it in list
        if 1.0 in sigma_list:
            ref_scale_idx = sigma_list.index(1.0)
        else:
            ref_scale_idx = 0
        S_reference = S_per_scale[ref_scale_idx]

        # --- Dynamic K: Significant Components ---
        # σ_i > 0.1 × σ_1 determines which components carry signal
        dynamic_k = self.determine_dynamic_k(
            S_reference,
            threshold_ratio=self.config.dynamic_k_threshold_ratio,
        )

        # --- U_basis: Top-K Right Singular Vectors ---
        # Re-compute SVD at reference scale to get full Vt for weighted distance
        ref_sigma = self._sanitize_sigma(
            sigma_list[ref_scale_idx], self.config.sigma_scale_epsilon
        )
        G_ref_scaled = torch.nan_to_num(
            G / ref_sigma, nan=0.0, posinf=0.0, neginf=0.0
        )
        _, _, Vt_ref = torch.linalg.svd(G_ref_scaled, full_matrices=False)
        # Gauge-fix each component consistently
        u_basis = Vt_ref  # [N, 8, H] — full basis, caller uses dynamic_k to slice

        # --- Directional Antagonism ---
        # Scale by the dominant singular value from reference scale
        antagonism = u_axis_fixed * S_reference[:, 0:1]

        # --- Dipole State ---
        n_valid_articles = dipole_valid.sum().item()
        if n_valid_articles > N / 2:
            dipole_state = "DIPOLE_STATE"
        else:
            dipole_state = "BLINKER_STATE"

        return SpectralPolarityResult(
            u_axis=u_axis_fixed.to(dtype=dtype),
            evr=evr_aggregated.to(dtype=dtype),
            dipole_valid=dipole_valid,
            probe_magnitudes=probe_magnitudes.to(dtype=dtype),
            singular_values=S_reference.to(dtype=dtype),
            antagonism=antagonism.to(dtype=dtype),
            n_persistent_scales=n_persistent,
            evr_per_scale=evr_stacked.to(dtype=dtype),
            dipole_state=dipole_state,
            dynamic_k=dynamic_k,
            u_basis=u_basis.to(dtype=dtype),
        )

    def compute_spectral_distance(
        self,
        cls_per_bot_list: List[torch.Tensor],
        u_axis: torch.Tensor,  # [N, H]
    ) -> torch.Tensor:
        """Compute Euclidean distance between poles along u_axis. Returns [N].

        DEPRECATED: Use compute_weighted_spectral_distance for signal-weighted subspace.
        """
        G = torch.stack(cls_per_bot_list, dim=0)  # [N, 8, H]
        projs = torch.bmm(G, u_axis.unsqueeze(-1)).squeeze(-1)  # [N, 8]
        proj_pos = projs.max(dim=1).values
        proj_neg = projs.min(dim=1).values
        return (proj_pos - proj_neg).abs()

    def determine_dynamic_k(
        self,
        singular_values: torch.Tensor,  # [N, 8]
        threshold_ratio: float = 0.1,
    ) -> torch.Tensor:
        """
        Determine dynamic K: number of significant spectral components per article.

        A component is significant if σ_i > threshold_ratio × σ_1 (the dominant).
        This avoids hardcoding k=3 and adapts to each article's spectral structure.

        Args:
            singular_values: [N, 8] singular value spectrum per article
            threshold_ratio: fraction of σ_1 below which components are noise (default 0.1)

        Returns:
            [N] tensor of integers: dynamic K per article (minimum 1)
        """
        # σ_1 is the largest singular value
        sigma_1 = singular_values[:, 0:1]  # [N, 1]
        threshold = threshold_ratio * sigma_1  # [N, 1]

        # Count how many components exceed the threshold
        significant = singular_values > threshold  # [N, 8] bool
        k_per_article = significant.sum(dim=1).clamp(min=1)  # [N], at least 1

        return k_per_article

    def compute_weighted_spectral_distance(
        self,
        emb_a: torch.Tensor,  # [N, H] or [H]
        emb_b: torch.Tensor,  # [N, H] or [H]
        singular_values: torch.Tensor,  # [N, 8] or [N, K]
        Vt: torch.Tensor,  # [N, 8, H] or [N, K, H] — right singular vectors
        dynamic_k: Optional[torch.Tensor] = None,  # [N] — if None, use all components
    ) -> torch.Tensor:
        """
        Compute signal-weighted spectral distance between two embeddings.

        d_spectral = sqrt( sum_i( σ_i × (proj_a_i - proj_b_i)^2 ) )

        This weights movement through the spectral subspace by the singular value
        "mass" of each component. High-variance directions (large σ) contribute
        more to the distance than low-variance directions.

        Args:
            emb_a: [N, H] first embedding (e.g., start position)
            emb_b: [N, H] second embedding (e.g., end position)
            singular_values: [N, K] singular values for weighting
            Vt: [N, K, H] right singular vectors defining the subspace
            dynamic_k: [N] optional per-article component count. If provided,
                       masks out components beyond K for each article.

        Returns:
            [N] weighted spectral distance per article
        """
        # Handle 1D input (single embedding)
        if emb_a.dim() == 1:
            emb_a = emb_a.unsqueeze(0)
        if emb_b.dim() == 1:
            emb_b = emb_b.unsqueeze(0)

        N = emb_a.shape[0]
        K = Vt.shape[1]  # Number of spectral components

        # Project both embeddings onto the spectral basis
        # proj = Vt @ emb.T → [N, K, H] @ [N, H, 1] → [N, K, 1] → [N, K]
        proj_a = torch.bmm(Vt, emb_a.unsqueeze(-1)).squeeze(-1)  # [N, K]
        proj_b = torch.bmm(Vt, emb_b.unsqueeze(-1)).squeeze(-1)  # [N, K]

        # Compute delta in spectral coordinates
        delta = proj_a - proj_b  # [N, K]

        # Weight by singular values: σ_i × (Δx_i)^2
        # Use first K singular values
        sigma = singular_values[:, :K]  # [N, K]
        sigma_term = self._sigma_weight_term(sigma, self.config.weighted_sigma_mode)
        weighted_sq = sigma_term * (delta ** 2)  # [N, K]

        # Apply dynamic K masking if provided
        if dynamic_k is not None:
            # Create mask: [N, K] where mask[n, k] = 1 if k < dynamic_k[n]
            k_range = torch.arange(K, device=singular_values.device).unsqueeze(0)  # [1, K]
            mask = k_range < dynamic_k.unsqueeze(1)  # [N, K]
            weighted_sq = weighted_sq * mask.float()

        # Sum and sqrt
        d_spectral = torch.sqrt(
            weighted_sq.sum(dim=1).clamp(min=self.config.normalization_epsilon)
        )  # [N]

        return d_spectral

    def compute_whitened_spectral_distance(
        self,
        emb_a: torch.Tensor,  # [N, H] or [H]
        emb_b: torch.Tensor,  # [N, H] or [H]
        singular_values: torch.Tensor,  # [N, 8] or [N, K]
        Vt: torch.Tensor,  # [N, 8, H] or [N, K, H] — right singular vectors
        dynamic_k: Optional[torch.Tensor] = None,  # [N] — if None, use all components
        epsilon: Optional[float] = None,  # Optional override floor for singular values
    ) -> torch.Tensor:
        """
        Compute whitened (inverse-weighted) spectral distance between two embeddings.

        d_whitened = sqrt( sum_i( (proj_a_i - proj_b_i)^2 / σ_i^2 ) )

        This AMPLIFIES movement in low-variance directions (small σ). Minor axes
        become hypersensitive — "lies in the details" are detected because small
        perturbations in low-variance directions produce large whitened distances.

        Use Case: Lie Detection
        - Weighted distance: measures "how far did you move in important directions?"
        - Whitened distance: measures "did you move suspiciously in unimportant directions?"

        A PHANTOM article will have small weighted distance (didn't move much in
        major axes) but large whitened distance (moved suspiciously in minor axes).

        Args:
            emb_a: [N, H] first embedding (e.g., start position)
            emb_b: [N, H] second embedding (e.g., end position)
            singular_values: [N, K] singular values for inverse weighting
            Vt: [N, K, H] right singular vectors defining the subspace
            dynamic_k: [N] optional per-article component count. If provided,
                       masks out components beyond K for each article.
            epsilon: floor for singular values to prevent division by zero

        Returns:
            [N] whitened spectral distance per article
        """
        # Handle 1D input (single embedding)
        if emb_a.dim() == 1:
            emb_a = emb_a.unsqueeze(0)
        if emb_b.dim() == 1:
            emb_b = emb_b.unsqueeze(0)

        N = emb_a.shape[0]
        K = Vt.shape[1]  # Number of spectral components

        # Project both embeddings onto the spectral basis
        proj_a = torch.bmm(Vt, emb_a.unsqueeze(-1)).squeeze(-1)  # [N, K]
        proj_b = torch.bmm(Vt, emb_b.unsqueeze(-1)).squeeze(-1)  # [N, K]

        # Compute delta in spectral coordinates
        delta = proj_a - proj_b  # [N, K]

        # Inverse weight by singular values: (Δx_i / σ_i)^2
        # Floor σ to prevent explosion on near-zero components
        sigma = singular_values[:, :K]  # [N, K]
        sigma_term = self._sigma_weight_term(sigma, self.config.whitened_sigma_mode)
        epsilon_floor = (
            self.config.whitened_sigma_epsilon if epsilon is None else float(epsilon)
        )
        sigma_term = sigma_term.clamp(min=epsilon_floor)
        whitened_sq = (delta / sigma_term) ** 2  # [N, K]

        # Apply dynamic K masking if provided
        if dynamic_k is not None:
            k_range = torch.arange(K, device=singular_values.device).unsqueeze(0)
            mask = k_range < dynamic_k.unsqueeze(1)  # [N, K]
            whitened_sq = whitened_sq * mask.float()

        # Sum and sqrt
        d_whitened = torch.sqrt(
            whitened_sq.sum(dim=1).clamp(min=self.config.normalization_epsilon)
        )  # [N]

        return d_whitened

    def compute_phantom_ratio(
        self,
        emb_a: torch.Tensor,  # [N, H]
        emb_b: torch.Tensor,  # [N, H]
        singular_values: torch.Tensor,  # [N, K]
        Vt: torch.Tensor,  # [N, K, H]
        dynamic_k: Optional[torch.Tensor] = None,
    ) -> torch.Tensor:
        """
        Compute phantom detection ratio: whitened / weighted distance.

        ratio = d_whitened / d_weighted

        Interpretation:
        - ratio ≈ 1: Movement is proportional across all axes (HONEST)
        - ratio >> 1: Movement concentrated in minor axes (PHANTOM - suspicious)
        - ratio << 1: Movement concentrated in major axes (TAUTOLOGY - obvious)

        This is the inverse of the Panic Function's divergence ratio.
        The Panic Function uses: Δ = W_actual / d_spectral (Walker / Spectral)
        This uses: d_whitened / d_weighted (minor-axis / major-axis sensitivity)

        Args:
            emb_a, emb_b: embeddings to compare
            singular_values, Vt: spectral decomposition
            dynamic_k: optional dynamic component count

        Returns:
            [N] phantom ratio per article (>1 indicates phantom behavior)
        """
        d_weighted = self.compute_weighted_spectral_distance(
            emb_a, emb_b, singular_values, Vt, dynamic_k
        )
        d_whitened = self.compute_whitened_spectral_distance(
            emb_a, emb_b, singular_values, Vt, dynamic_k
        )

        # Avoid division by zero
        ratio = d_whitened / d_weighted.clamp(min=self.config.normalization_epsilon)

        return ratio

    def compute_expected_spectral_distance(
        self,
        singular_values: torch.Tensor,  # [N, 8]
        dynamic_k: Optional[torch.Tensor] = None,  # [N]
    ) -> torch.Tensor:
        """
        Compute expected spectral distance for random unit displacement.

        E[d] = sqrt( sum_i(σ_i) / K ) — the "null model" distance

        This represents the expected distance if we moved uniformly in all
        spectral directions. Used as denominator in the Panic Function.

        Args:
            singular_values: [N, 8] singular value spectrum
            dynamic_k: [N] optional dynamic component count

        Returns:
            [N] expected spectral distance per article
        """
        N = singular_values.shape[0]
        K = singular_values.shape[1]

        if dynamic_k is None:
            # Use all components
            sigma_sum = singular_values.sum(dim=1)  # [N]
            k_effective = torch.full((N,), K, device=singular_values.device, dtype=torch.float32)
        else:
            # Mask to dynamic K
            k_range = torch.arange(K, device=singular_values.device).unsqueeze(0)  # [1, K]
            mask = k_range < dynamic_k.unsqueeze(1)  # [N, K]
            sigma_sum = (singular_values * mask.float()).sum(dim=1)  # [N]
            k_effective = dynamic_k.float().clamp(min=1)

        # E[d] = sqrt(sum(σ) / K)
        expected = torch.sqrt(sigma_sum / k_effective)

        return expected

    def _compute_single_scale_legacy(
        self,
        G: torch.Tensor,
        dtype: torch.dtype,
    ) -> SpectralPolarityResult:
        """
        Legacy single-scale computation (no multi-scale sweep).
        Kept for backward compatibility and ablation studies.
        """
        N, n_probes, H = G.shape

        # SVD at scale=1.0
        U, S, Vt = torch.linalg.svd(G, full_matrices=False)
        u_axis = Vt[:, 0, :]

        S_sq = S ** 2
        total_var = S_sq.sum(dim=1).clamp(min=self.config.normalization_epsilon)
        evr = S_sq[:, 0] / total_var

        dipole_valid = evr >= self.config.evr_threshold

        # Gauge fixing
        u_axis = self._gauge_fix(u_axis, G, self.config.reference_probe_idx)

        # Probe magnitudes
        probe_magnitudes = torch.bmm(G, u_axis.unsqueeze(2)).squeeze(2)

        # Antagonism
        antagonism = u_axis * S[:, 0:1]

        # Dynamic K
        dynamic_k = self.determine_dynamic_k(
            S,
            threshold_ratio=self.config.dynamic_k_threshold_ratio,
        )

        # Fake multi-scale outputs for compatibility
        n_persistent = dipole_valid.long()
        evr_per_scale = evr.unsqueeze(1)

        n_valid = dipole_valid.sum().item()
        dipole_state = "DIPOLE_STATE" if n_valid > N / 2 else "BLINKER_STATE"

        return SpectralPolarityResult(
            u_axis=u_axis.to(dtype=dtype),
            evr=evr.to(dtype=dtype),
            dipole_valid=dipole_valid,
            probe_magnitudes=probe_magnitudes.to(dtype=dtype),
            singular_values=S.to(dtype=dtype),
            antagonism=antagonism.to(dtype=dtype),
            n_persistent_scales=n_persistent,
            evr_per_scale=evr_per_scale.to(dtype=dtype),
            dipole_state=dipole_state,
            dynamic_k=dynamic_k,
            u_basis=Vt.to(dtype=dtype),
        )

\end{CodeBlock}
\end{document}
